	\section{Estado del Arte}
	La transformación digital ha impulsado el desarrollo de plataformas orientadas a la gestión
	electrónica de denuncias, quejas y mecanismos de protección de derechos. Organizaciones
	gubernamentales, corporativas y académicas han incorporado sistemas capaces de automatizar la
	recepción de reportes, el seguimiento de expedientes y la administración de evidencia electrónica.
	
	Esa evolución, sin embargo, no ha sido homogénea. Algunos sistemas privilegian la validez jurídica y
	el cumplimiento normativo, mientras que otros se concentran en la anonimización de datos, la
	movilidad, la trazabilidad o la incorporación de mecanismos de seguridad informática. El resultado es
	que buena parte de las soluciones disponibles resulta insuficiente para contextos institucionales
	complejos, en particular para entornos universitarios donde convergen aspectos administrativos,
	jurídicos, académicos y de protección de derechos humanos.
	
	En el caso de la DDP-IPN, una plataforma especializada debe atender exigencias que rebasan las de un
	sistema convencional de recepción de denuncias. Los aspectos críticos identificados son los
	siguientes:
	
	\begin{itemize}
		\item Protección e integridad de la información sensible.
		\item Automatización de procesos administrativos.
		\item Escalabilidad de la infraestructura tecnológica.
		\item Trazabilidad de expedientes digitales.
		\item Protección de identidad y anonimización del denunciante.
		\item Integración institucional entre distintas áreas administrativas.
		\item Adaptabilidad a procesos universitarios complejos.
	\end{itemize}
	
	Con ese marco, este capítulo analiza diversas plataformas nacionales e internacionales de denuncia
	digital y gestión institucional, con el objetivo de identificar las capacidades tecnológicas
	existentes, las limitaciones arquitectónicas predominantes y las brechas funcionales que justifican
	el desarrollo de una solución especializada para el Instituto Politécnico Nacional.
	
	\subsection{Criterios de evaluación}
	Para analizar las plataformas seleccionadas bajo una misma óptica, se definieron criterios de
	evaluación que abarcan tanto aspectos funcionales como arquitectónicos:
	
	\begin{itemize}
		\item \textbf{Escalabilidad arquitectónica:} capacidad del sistema para absorber crecimiento en
		usuarios, servicios y procesamiento sin degradar de forma significativa su rendimiento.
	
		\item \textbf{Seguridad y confidencialidad:} mecanismos con que protege la información sensible,
		entre ellos cifrado, autenticación y control de acceso.
	
		\item \textbf{Anonimización y protección de identidad:} nivel de protección que ofrece al
		denunciante para preservar su identidad y confidencialidad.
	
		\item \textbf{Trazabilidad y seguimiento:} posibilidad de monitorear el estado y el avance de los
		expedientes o denuncias digitales.
	
		\item \textbf{Automatización administrativa:} uso de tecnologías que reducen el trabajo manual
		mediante reglas de negocio, clasificación automática o asistencia inteligente.
	
		\item \textbf{Soporte multimedia:} capacidad de almacenar y gestionar evidencia digital como
		imágenes, documentos, audio y video.
	
		\item \textbf{Interoperabilidad institucional:} facilidad para integrarse con otros servicios,
		sistemas o plataformas de la propia institución.
	
		\item \textbf{Adaptación al entorno universitario:} grado de adecuación funcional y administrativa
		a instituciones educativas públicas.
	\end{itemize}
	
	Sobre estos criterios se construyó el análisis comparativo que se presenta a continuación.
	
	\subsection{Plataformas gubernamentales}
	Esta sección revisa las herramientas tecnológicas que el Estado mexicano ha implementado para
	recibir, gestionar y dar seguimiento a denuncias ciudadanas. A partir de soluciones federales, como
	el Sistema Integral de Denuncias Ciudadanas (SIDEC), y de iniciativas locales, como la plataforma
	Denuncia Digital de la Ciudad de México, se examinan sus capacidades funcionales, su arquitectura y
	los retos que aún enfrentan en materia de interoperabilidad, automatización y descentralización.
	
	\subsubsection{Sistema Integral de Denuncias Ciudadanas (SIDEC)}
	El SIDEC es una de las principales plataformas gubernamentales de recepción de quejas
	administrativas en México y centraliza las denuncias relacionadas con servidores públicos federales
	\cite{sfp_sidec}.
	
	En el plano funcional, el sistema recibe denuncias de forma continua mediante formularios
	electrónicos y ofrece mecanismos básicos de seguimiento a través de folios únicos y contraseñas de
	acceso \cite{amexcid_sidec}. Admite además evidencia multimedia en formato de documento, audio y
	video. Su relevancia como referente se refuerza con los trabajos de mejora que la Secretaría de la
	Función Pública ha impulsado junto con el BID y GovLab \cite{sfp_bid_govlab_sidec}.
	
	El análisis arquitectónico, en cambio, revela varias limitaciones técnicas:
	
	\begin{itemize}
		\item Arquitectura predominantemente centralizada.
		\item Procesos administrativos parcialmente manuales.
		\item Seguimiento limitado para el denunciante.
		\item Ausencia de automatización inteligente.
		\item Escasa modularidad de servicios.
	\end{itemize}
	
	\subsubsection{Denuncia Digital y sistemas estatales}
	Diversas entidades federativas han incorporado plataformas digitales para recibir denuncias y
	querellas. El caso más representativo es el sistema Denuncia Digital de la Ciudad de México
	\cite{cdmx_denuncia_digital}, que permite iniciar procedimientos por delitos no violentos mediante
	autenticación electrónica con la Llave CDMX o con firma electrónica gubernamental.
	
	La comparación entre entidades revela una fragmentación tecnológica considerable
	\cite{impunidadcero_denuncia_digital}: algunas plataformas ya admiten videodenuncias y ratificación
	remota, mientras otras siguen operando con modelos híbridos o trámites presenciales. Aunque
	representan un avance real en digitalización administrativa, estas soluciones se quedan cortas en
	interoperabilidad, automatización y protección avanzada de identidad.
	
	\subsubsection{Servicios de procuración digital: MP Virtual y denuncia anónima}
	La Fiscalía General de Justicia de la Ciudad de México ha sumado servicios digitales
	complementarios, entre ellos el MP Virtual y el sistema de denuncia anónima
	\cite{fgjcdmx_servicios_linea}.
	
	El módulo MP Virtual formaliza a distancia determinados procedimientos jurídicos, lo que facilita la
	validación documental y reduce la necesidad de comparecencia presencial \cite{fgjcdmx_mp_virtual}.
	El sistema de denuncia anónima, por su parte, implementa mecanismos básicos de protección de
	identidad para reportes de delitos de alto impacto.
	
	Ambos servicios, no obstante, siguen apoyados en infraestructuras centralizadas y procesos
	administrativos tradicionales, y carecen de componentes de inteligencia artificial, clasificación
	automática y analítica institucional.
	
	\subsection{Aplicaciones móviles y plataformas ciudadanas}
	Junto a las plataformas web institucionales ha surgido una generación de aplicaciones móviles y
	canales ciudadanos que privilegian la accesibilidad y la participación colaborativa por encima de la
	formalidad administrativa.
	
	\subsubsection{Aplicaciones de denuncia georreferenciada}
	La evolución de los servicios digitales ha impulsado aplicaciones móviles orientadas a la denuncia
	ciudadana y la participación colaborativa. Uno de los referentes más conocidos es la aplicación
	Incorruptible, diseñada para registrar actos de corrupción mediante geolocalización
	\cite{nacion321_apps_corrupcion}. Sus funcionalidades principales son:
	
	\begin{itemize}
		\item Registro georreferenciado de incidentes.
		\item Publicación colaborativa de reportes.
		\item Integración de evidencia multimedia.
		\item Cuantificación económica de los daños reportados.
	\end{itemize}
	
	En una línea complementaria, el Sistema de Atención Mexiquense (SAM) integra canales web, móviles y
	telefónicos para canalizar denuncias hacia distintas dependencias administrativas
	\cite{secogem_sam}.
	
	Su accesibilidad y facilidad de uso son evidentes, pero este tipo de aplicaciones queda limitado en
	trazabilidad jurídica, validación institucional y protección de datos sensibles.
	
	\subsection{Plataformas internacionales}
	El mercado internacional de plataformas de denuncia y cumplimiento normativo ofrece soluciones
	maduras, dirigidas sobre todo a sectores corporativos y regulados. Destacan por sus mecanismos de
	seguridad, anonimización y administración de casos, aunque suelen tropezar con costos de
	implementación elevados, dependencia de proveedores externos y dificultades para adaptarse a marcos
	institucionales específicos.
	
	\subsubsection{Whistlelink}
	Whistlelink gestiona denuncias corporativas de forma anónima y destaca por su soporte multilingüe y
	por sus canales bidireccionales de comunicación anónima \cite{whistlelink}. Su mayor virtud es
	reducir las barreras de acceso para denunciantes de distintos países; en contraparte, ofrece poca
	personalización institucional y se adapta con dificultad a procesos administrativos complejos.
	
	\subsubsection{EQS Integrity Line}
	EQS Integrity Line es una de las soluciones corporativas más robustas del mercado europeo, diseñada
	para cumplir con la normativa internacional de protección al denunciante \cite{eqs_integrity_line}.
	Incorpora herramientas avanzadas de administración de casos, monitoreo organizacional y generación
	de reportes. Su complejidad operativa y la carga administrativa que exige, sin embargo, complican su
	implementación en instituciones educativas públicas.
	
	\subsubsection{NAVEX / WhistleB}
	NAVEX y su solución WhistleB se apoyan en infraestructura en la nube para garantizar alta
	disponibilidad y escalabilidad \cite{navex_whistleb}. Entre sus fortalezas figuran:
	
	\begin{itemize}
		\item Integración con plataformas GRC.
		\item Infraestructura escalable.
		\item Gestión centralizada de cumplimiento.
		\item Administración avanzada de casos.
	\end{itemize}
	
	La dependencia de proveedores externos, en cambio, plantea restricciones de soberanía de datos que
	chocan con las políticas de las instituciones gubernamentales.
	
	\subsubsection{Whistleblower Software}
	Whistleblower Software sobresale por sus mecanismos de cifrado y por adoptar un modelo de
	arquitectura \textit{Zero-Knowledge} \cite{whistleblower_software}, enfoque que refuerza de manera
	notable la confidencialidad de la información y la protección del denunciante. Su orientación al
	entorno corporativo, no obstante, limita su encaje directo en estructuras universitarias públicas.
	
	\subsection{Plataformas académicas e institucionales}
	Esta sección examina las herramientas e iniciativas digitales que instituciones educativas y
	organismos públicos han desarrollado para atender incidencias en el ámbito académico. Mediante el
	análisis de BullyButton, Ventanilla Escolar, Alerta IPN y el sistema de Denuncia en Línea de la
	UNAM se evalúan sus alcances operativos y protocolos de respuesta, y se identifican las brechas
	tecnológicas comunes a estos entornos: sobre todo la gestión jurídica de expedientes, la
	automatización del triaje y la falta de mecanismos avanzados para proteger la identidad de los
	usuarios.
	
	\subsubsection{BullyButton}
	BullyButton se enfoca en la detección y el reporte de casos de acoso escolar \cite{bully_button}. Su
	diseño privilegia la simplicidad operativa y la facilidad de uso para estudiantes, pero carece de
	mecanismos formales de trazabilidad jurídica y de administración avanzada de expedientes.
	
	\subsubsection{Ventanilla Escolar}
	Ventanilla Escolar opera como canal de orientación y comunicación entre las instituciones educativas
	y las autoridades federales \cite{sep_ventanilla_escolar}. Centraliza información relevante para la
	comunidad educativa, aunque no cubre el seguimiento de denuncias ni la administración formal de
	expedientes.
	
	\subsubsection{Alerta IPN}
	Alerta IPN es uno de los principales mecanismos institucionales de reporte dentro del Instituto
	Politécnico Nacional \cite{ipn_alerta} y se orienta sobre todo a emergencias y situaciones de riesgo
	dentro de las instalaciones. Sus limitaciones se concentran en:
	
	\begin{itemize}
		\item Gestión jurídica de expedientes.
		\item Clasificación administrativa automatizada.
		\item Seguimiento integral de casos.
		\item Protección avanzada de identidad.
		\item Automatización inteligente de procesos.
	\end{itemize}
	
	\subsubsection{UNAM: Denuncia en Línea}
	La plataforma de Denuncia en Línea de la UNAM es uno de los referentes más relevantes de gestión
	universitaria de denuncias en México \cite{unam_ddu_denuncia}. Su mayor fortaleza son los protocolos
	institucionales consolidados y los procedimientos administrativos claramente definidos. El análisis
	tecnológico, en cambio, evidencia limitaciones en:
	
	\begin{itemize}
		\item Escalabilidad arquitectónica.
		\item Integración de inteligencia artificial.
		\item Automatización del triaje.
		\item Modularidad de servicios.
		\item Gestión inteligente de expedientes.
	\end{itemize}
	
	\subsection{Tabla comparativa general}
	Con base en los criterios definidos al inicio del capítulo, se comparó el conjunto de plataformas
	analizadas para contrastar sus capacidades funcionales, sus limitaciones arquitectónicas y su nivel
	de adaptación a entornos institucionales complejos. La Tabla~\ref{tab:comparativa_plataformas}
	resume las características identificadas durante la investigación.
	
	\rowcolors{2}{rowblue}{white}
	
	\footnotesize
	\renewcommand{\arraystretch}{1.2}
	\setlength{\tabcolsep}{3pt}
	
	\begin{longtable}{|p{2.7cm}|p{1.55cm}|p{1.7cm}|p{1.5cm}|p{1.4cm}|p{1.7cm}|p{1.5cm}|p{1.8cm}|}
		\caption{Comparativa general de plataformas de denuncia digital}
		\label{tab:comparativa_plataformas} \\
		\hline
		\rowcolor{headerblue}
		\headercell{Plataforma} &
		\headercell{Seguridad} &
		\headercell{Anonimato} &
		\headercell{Segui\-miento} &
		\headercell{Multi\-media} &
		\headercell{Automati\-zación} &
		\headercell{Escala\-bilidad} &
		\headercell{Adaptación universitaria} \\
		\hline
		\endfirsthead
	
		\hline
		\rowcolor{headerblue}
		\headercell{Plataforma} &
		\headercell{Seguridad} &
		\headercell{Anonimato} &
		\headercell{Segui\-miento} &
		\headercell{Multi\-media} &
		\headercell{Automati\-zación} &
		\headercell{Escala\-bilidad} &
		\headercell{Adaptación universitaria} \\
		\hline
		\endhead
	
		SIDEC & Media & Baja & Media & Sí & Baja & Media & No \\ \hline
		Denuncia Digital CDMX & Media & Baja & Media & Sí & Baja & Media & No \\ \hline
		MP Virtual & Media & Media & Media & Sí & Baja & Media & No \\ \hline
		Incorruptible & Baja & Media & Baja & Sí & Baja & Alta & No \\ \hline
		SAM & Media & Baja & Media & Sí & Baja & Media & No \\ \hline
		Whistlelink & Alta & Alta & Alta & Sí & Media & Alta & Parcial \\ \hline
		EQS Integrity Line & Alta & Alta & Alta & Sí & Alta & Alta & Parcial \\ \hline
		NAVEX / WhistleB & Alta & Alta & Alta & Sí & Alta & Alta & Parcial \\ \hline
		Whistleblower Software & Muy alta & Muy alta & Alta & Sí & Alta & Alta & Parcial \\ \hline
		BullyButton & Media & Media & Baja & Limitado & Baja & Media & Sí \\ \hline
		Ventanilla Escolar & Media & Baja & Baja & No & Baja & Media & Sí \\ \hline
		Alerta IPN & Media & Baja & Media & Sí & Baja & Media & Sí \\ \hline
		UNAM Denuncia en Línea & Media & Media & Alta & Sí & Media & Media & Sí \\ \hline
	\end{longtable}
	
	\normalsize
	\setlength{\tabcolsep}{6pt}
	\renewcommand{\arraystretch}{1}
	
	\subsection{Brechas tecnológicas identificadas}
	El análisis comparativo muestra que, si bien existen múltiples plataformas dedicadas a la gestión
	digital de denuncias y quejas, ninguna reúne por completo los requerimientos tecnológicos,
	administrativos y de protección institucional que exige un entorno universitario público como el
	Instituto Politécnico Nacional. Las brechas principales son:
	
	\begin{itemize}
		\item Dependencia de arquitecturas centralizadas tradicionales.
		\item Escasa automatización inteligente de los procesos administrativos.
		\item Mecanismos limitados de anonimización y protección de identidad.
		\item Baja interoperabilidad entre servicios institucionales.
		\item Deficiencias en la trazabilidad integral de los expedientes digitales.
		\item Escasa adaptación a estructuras universitarias complejas.
		\item Ausencia de arquitecturas híbridas escalables pensadas para instituciones académicas
		públicas.
	\end{itemize}
	
	Estas limitaciones justifican el desarrollo de una plataforma especializada que combine seguridad
	informática avanzada, modularidad arquitectónica, automatización administrativa y adaptación
	específica al contexto institucional de la DDP-IPN.
